\chapter{孤立奇点与留数}
\label{chap:isolated-singularities-residues}

Taylor 级数描述全纯函数在普通点附近的行为。若函数只在删去中心的圆盘上
全纯，正幂级数不再足够：负幂记录函数接近缺失点时的发散和振荡。Laurent
展开把这种局部行为编码为双向级数，其中 \((z-a)^{-1}\) 的系数又恰好决定
闭曲线积分。留数法的计算效率来自这一事实，而不是一种脱离 Cauchy 定理的
形式规则。

\section{环域上的 Laurent 展开}

设

\[
 A(a;r,R)=\{z\in\C:r<|z-a|<R\},
 \qquad 0\le r<R\le\infty .
\]

双向级数 \(\sum_{n=-\infty}^{\infty}a_n(z-a)^n\) 的含义是正幂部分
\(\sum_{n\ge0}\) 与负幂部分 \(\sum_{n\ge1}a_{-n}(z-a)^{-n}\) 分别收敛。

\begin{theorem}[Laurent 展开定理]
\label{thm:laurent-expansion}
若 \(f\) 在环域 \(A(a;r,R)\) 全纯，则存在唯一系数列
\((a_n)_{n\in\Z}\)，使

\begin{equation}
 f(z)=\sum_{n=-\infty}^{\infty}a_n(z-a)^n
 \label{eq:laurent-series}
\end{equation}

在环域内成立，并在每个紧子环域上一致绝对收敛。对任意
\(r<\rho<R\)，系数由

\begin{equation}
 a_n=\frac1{2\pi i}\int_{|\zeta-a|=\rho}
 \frac{f(\zeta)}{(\zeta-a)^{n+1}}\,d\zeta,
 \qquad n\in\Z,
 \label{eq:laurent-coefficients}
\end{equation}

给出，且右端与 \(\rho\) 无关。
\end{theorem}

\begin{proof}
固定 \(z\in A(a;r,R)\)，选择

\[
 r<\rho<|z-a|<\sigma<R.
\]

记两个正向圆周之差为 \(\Gamma=C_\sigma-C_\rho\)。它在环域外的绕数为零，
在 \(z\) 处的绕数为一；单独一个圆周则不满足前一个条件。
在 \(z\) 周围取充分小的正向圆周 \(C_\eta(z)\)，使其闭圆盘留在两圆之间。
链 \(\Gamma-C_\eta(z)\) 在删去 \(z\) 的环域外绕数全为零，故可对
\(f(\zeta)/(\zeta-z)\) 使用推论~\ref{cor:cauchy-cycle-form}。
令 \(\eta\downarrow0\)，小圆积分趋于 \(2\pi i f(z)\)。
由此得到差链的 Cauchy 公式，其中内边界取反向：

\[
 f(z)=\frac1{2\pi i}\int_{|\zeta-a|=\sigma}
 \frac{f(\zeta)}{\zeta-z}\,d\zeta
 +
 \frac1{2\pi i}\int_{|\zeta-a|=\rho}
 \frac{f(\zeta)}{z-\zeta}\,d\zeta.
\]

在外圆上有

\[
 \frac1{\zeta-z}
 =\sum_{n=0}^{\infty}\frac{(z-a)^n}{(\zeta-a)^{n+1}},
\]

在内圆上则有

\[
 \frac1{z-\zeta}
 =\sum_{k=0}^{\infty}\frac{(\zeta-a)^k}{(z-a)^{k+1}}.
\]

两级数分别在积分圆周上一致收敛，故可逐项积分；所得系数正是
式\eqref{eq:laurent-coefficients}中 \(n\ge0\) 和 \(n=-k-1\) 的情形。
若改换积分半径，两圆之间的被积函数
\(f(\zeta)(\zeta-a)^{-n-1}\) 全纯，Cauchy 定理说明两积分相同。

对任意紧子环域
\(r<r_1\le|z-a|\le R_1<R\)，可固定
\(r<\rho<r_1\le R_1<\sigma<R\)。上面的两个几何级数在该紧子环域上各有
统一比值小于一，因而所得 Laurent 级数一致绝对收敛。

若另一个 Laurent 级数也表示 \(f\)，在任意中间圆周上一致收敛允许逐项
积分。乘以 \((z-a)^{-m-1}\) 后，只有指数为 \(-1\) 的项积分不为零，
由此恢复 \(a_m\)，故系数唯一。
\end{proof}

系数公式还给出双向的 Cauchy 估计。若
\(M_\rho=\max_{|\zeta-a|=\rho}|f(\zeta)|\)，则对每个 \(n\in\Z\)，
\begin{equation}
 |a_n|\le M_\rho\rho^{-n}.
 \label{eq:laurent-coefficient-estimate}
\end{equation}
当 \(n<0\) 时，右端含有正幂 \(\rho^{-n}\)；这正是让圆周缩向中心后消去
负次系数的工具。

负幂部分

\[
 \sum_{n=1}^{\infty}a_{-n}(z-a)^{-n}
\]

称为 Laurent 展开的\term{主部}。同一个函数在不同环域上的展开可能不同。
例如

\[
 \frac1{z(1-z)}
\]

在 \(0<|z|<1\) 中展开为
\(\sum_{n\ge0}z^{n-1}\)，而在 \(|z|>1\) 中展开为
\(-\sum_{n\ge0}z^{-n-2}\)。写 Laurent 级数时必须同时说明环域。

\section{可去奇点、极点和本性奇点}

若 \(f\) 在某个删心圆盘 \(0<|z-a|<R\) 全纯，则称 \(a\) 是 \(f\) 的
\term{孤立奇点}。Laurent 主部给出三种互斥情形。

\begin{theorem}[孤立奇点分类]
\label{thm:isolated-singularity-classification}
设 \(a\) 是 \(f\) 的孤立奇点。

\begin{enumerate}[label=(\roman*)]
 \item 若所有负次系数都为零，则 \(a\) 是可去奇点，\(f\) 可唯一延拓为
       \(a\) 邻域的全纯函数。
 \item 若主部只有有限项，且 \(a_{-m}\ne0\) 是最低次项，则 \(a\) 是
       \(m\) 阶极点，并有
       \(f(z)=(z-a)^{-m}g(z)\)，其中 \(g\) 在 \(a\) 附近全纯且
       \(g(a)\ne0\)。
 \item 若主部有无穷多非零项，则 \(a\) 是本性奇点。
\end{enumerate}
\end{theorem}

第一种情形直接把正幂级数在 \(a\) 处取值 \(a_0\)。第二种情形提出最低负幂
即可。由于 Laurent 展开唯一，不存在第四种情形。

三类情形可由熟悉的函数看清。函数
\[
 \frac{\sin(z-a)}{z-a}
\]
在 \(a\) 可去，延拓值为 \(1\)；\(1/\sin(z-a)\) 在 \(a\) 有单极点；
\(e^{1/(z-a)}\) 则有无穷多个负次项，所以是本性奇点。最后一个函数沿
\(z=a-t^{-1}\)、\(t\to\infty\) 趋于零，沿 \(z=a+t^{-1}\) 却趋于无穷。
因而沿一条路径有界甚至存在极限，都不足以判定可去性；可去奇点定理要求整个
删心邻域上的控制。

\begin{theorem}[Riemann 可去奇点定理]
\label{thm:riemann-removable-singularity}
设 \(f\) 在 \(0<|z-a|<R\) 全纯。以下条件等价：

\begin{enumerate}[label=(\roman*)]
 \item \(a\) 是可去奇点；
 \item \(f\) 在 \(a\) 的某个删心邻域有界；
 \item \(\lim_{z\to a}f(z)\) 存在且有限。
\end{enumerate}
\end{theorem}

\begin{proof}
\((i)\Rightarrow(iii)\Rightarrow(ii)\) 来自延拓函数的连续性。设 (ii) 成立。
则存在 \(r_0\in(0,R)\) 和一个与半径无关的常数 \(M\)，使
\[
 |f(z)|\le M\qquad(0<|z-a|<r_0).
\]
对任意 \(0<\rho<r_0\) 和 \(n\ge1\)，Laurent 系数公式给

\[
 |a_{-n}|
 =\left|\frac1{2\pi i}\int_{|\zeta-a|=\rho}
 f(\zeta)(\zeta-a)^{n-1}\,d\zeta\right|
 \le M\rho^n.
\]

该估计对任意充分小的 \(\rho\) 成立，令 \(\rho\downarrow0\) 得
\(a_{-n}=0\)。主部消失，所以奇点可去。
\end{proof}

极点也有不依赖级数的判据：\(a\) 是极点当且仅当
\(|f(z)|\to\infty\) 当 \(z\to a\)。正向由分解
\(f=(z-a)^{-m}g\) 得到；反向令 \(1/f\) 在删心邻域定义，它趋于零并可去，
延拓后在 \(a\) 有某个正阶零点，取倒数便得到极点。

本性奇点的行为不能归结为有限阶发散。

\begin{theorem}[Casorati--Weierstrass 定理]
\label{thm:casorati-weierstrass}
若 \(a\) 是 \(f\) 的本性奇点，则任意删心邻域的像在 \(\C\) 中稠密。
\end{theorem}

\begin{proof}
若某个删心邻域的像不稠密，则存在 \(w_0\in\C\) 和 \(\varepsilon>0\)，
使 \(|f(z)-w_0|\ge\varepsilon\)。于是

\[
 g(z)=\frac1{f(z)-w_0}
\]

有界，按可去奇点定理延拓到 \(a\)。若 \(g(a)\ne0\)，则
\(f=w_0+1/g\) 可去；若 \(g(a)=0\)，则 \(1/g\) 有极点，\(f\) 也有极点。
两者都与本性奇点矛盾。
\end{proof}

稠密不等于每个值都在每个删心邻域取得。更强的 Picard 大定理只允许至多
遗漏一个值，但其证明不在当前理论链中。

若一个函数在域 \(D\) 中每点附近或者全纯，或者只有以该点为极点的孤立奇点，
就称它在 \(D\) 上\term{亚纯}；可去点在这里已按全纯方式补上。
极点处也可把函数值记作 \(\infty\)，但本书不把恒为 \(\infty\) 的对象算作亚纯函数。
这个局部条件保证极点在 \(D\) 内没有聚点。

亚纯函数在每点附近都可写成两个全纯函数之商，分母不恒为零：
在极点 \(a\) 附近用 \(g(z)/(z-a)^m\)，在普通点用 \(f/1\)。
反过来，全纯商的分母在一个零点附近可提出有限次幂；若分子也有零点，
约去公共幂次，便只剩全纯函数或有限阶极点。分子局部恒为零时，
所得商补成零函数。这说明两种局部描述等价。

\section{留数和留数定理}

设 \(a\) 是孤立奇点。Laurent 系数 \(a_{-1}\) 称为 \(f\) 在 \(a\) 的
\term{留数}，记为

\[
 \operatorname{Res}(f,a)=a_{-1}
 =\frac1{2\pi i}\int_{|z-a|=\rho}f(z)\,dz.
\]

若 \(a\) 是可去奇点，留数为零；反之留数为零并不能推出奇点可去，例如
\(1/(z-a)^2\)。若 \(a\) 是单极点，则

\[
 \operatorname{Res}(f,a)=\lim_{z\to a}(z-a)f(z).
\]

若 \(f=\phi/\psi\)，其中 \(\phi,\psi\) 在 \(a\) 附近全纯，
\(\psi(a)=0\)、\(\psi'(a)\ne0\)，则

\[
 \operatorname{Res}(f,a)=\frac{\phi(a)}{\psi'(a)}.
\]

对于 \(m\) 阶极点，

\begin{equation}
 \operatorname{Res}(f,a)
 =\frac1{(m-1)!}
 \left[\frac{d^{m-1}}{dz^{m-1}}\bigl((z-a)^m f(z)\bigr)\right]_{z=a}.
 \label{eq:residue-pole-formula}
\end{equation}

\begin{theorem}[留数定理]
\label{thm:residue-theorem}
设 \(D\) 为域，\(a_1,\ldots,a_N\in D\) 两两不同，函数 \(f\) 在
\(D\setminus\{a_1,\ldots,a_N\}\) 全纯。设闭分片 \(C^1\) 曲线
\(\gamma\subset D\setminus\{a_1,\ldots,a_N\}\) 满足
\(\operatorname{Ind}(\gamma,w)=0\) 对所有 \(w\notin D\) 成立。则

\begin{equation}
 \int_\gamma f(z)\,dz
 =2\pi i\sum_{j=1}^N
 \operatorname{Ind}(\gamma,a_j)\operatorname{Res}(f,a_j).
 \label{eq:residue-theorem}
\end{equation}
\end{theorem}

\begin{proof}
在每个 \(a_j\) 周围取互不相交、避开 \(\gamma\) 且闭含于 \(D\) 的小闭圆盘，
其正向边界记为 \(C_j\)。令
\[
 n_j=\operatorname{Ind}(\gamma,a_j),\qquad
 \Gamma=\gamma-\sum_{j=1}^N n_jC_j,\qquad
 \Omega=D\setminus\{a_1,\ldots,a_N\}.
\]
平面域删去有限多个点后仍为域（逐次使用上一章的删点绕行论证），曲线链
\(\Gamma\) 位于 \(\Omega\)。若 \(w\notin D\)，则题设给
\(\operatorname{Ind}(\gamma,w)=0\)，而每个小圆盘闭含于 \(D\)，故
\(\operatorname{Ind}(C_j,w)=0\)。若 \(w=a_k\)，小圆盘互不相交给出
\[
 \operatorname{Ind}(C_j,a_k)=
 \begin{cases}1,&j=k,\\0,&j\ne k,
 \end{cases}
\]
因而 \(\operatorname{Ind}(\Gamma,a_k)=n_k-n_k=0\)。这已经核对了
\(\operatorname{Ind}(\Gamma,w)=0\) 对所有 \(w\notin\Omega\) 成立。

函数 \(f\) 在 \(\Omega\) 全纯，故曲线链形式的 Cauchy 定理
（推论~\ref{cor:cauchy-cycle-form}）给出 \(\int_\Gamma f\,dz=0\)，即

\[
 \int_\gamma f(z)\,dz
 =\sum_{j=1}^N\operatorname{Ind}(\gamma,a_j)\int_{C_j}f(z)\,dz.
\]

Laurent 级数在每个 \(C_j\) 上一致收敛，逐项积分时只有
\((z-a_j)^{-1}\) 项留下 \(2\pi i\)，所以
\(\int_{C_j}f\,dz=2\pi i\operatorname{Res}(f,a_j)\)。代入上式即得
式\eqref{eq:residue-theorem}。
\end{proof}

若 \(\gamma\) 是简单闭曲线的正向边界，所有内部点绕数为 \(1\)、外部点为
\(0\)，公式退化为“边界积分等于内部留数之和的 \(2\pi i\) 倍”。奇点落在
曲线上时公式没有定义；必须改变路径、取主值，或另作局部处理。

\section{无穷远点处的奇点}

若 \(f\) 在 \(|z|>R\) 全纯，则令

\[
 F(w)=f(1/w).
\]

按 \(F\) 在 \(w=0\) 的类型，把 \(\infty\) 分为可去奇点、极点或本性奇点。
例如非常数多项式在无穷远有极点，常数函数在无穷远可去；
非多项式整函数在无穷远有本性奇点，而
\(1/z\) 在无穷远可去。

无穷远点的留数定义为

\begin{equation}
 \operatorname{Res}(f,\infty)
 =-\operatorname{Res}\left(\frac{f(1/w)}{w^2},0\right).
 \label{eq:residue-infinity}
\end{equation}

在局部坐标 \(w=1/z\) 中，
\[
 dz=-w^{-2}\,dw,\qquad
 f(z)\,dz=-\frac{f(1/w)}{w^2}\,dw.
\]
定义中的负号来自这个微分变换。此外，正向大圆在 \(w\) 平面中对应顺时针
小圆；计算积分时须同时处理微分和圆周方向。这并不表示全纯坐标变换
反转实平面的定向。若 \(f\) 在某个外圆域的
Laurent 展开写成
\[
 f(z)=\sum_{n=-\infty}^{\infty}c_nz^n,
\]
则代入定义可见 \(\operatorname{Res}(f,\infty)=-c_{-1}\)。若 \(f\) 在充分大的圆外
全纯，则对每个充分大的 \(R\) 都有

\[
 \operatorname{Res}(f,\infty)
 =-\frac1{2\pi i}\int_{|z|=R}f(z)\,dz.
\]

若 \(f\) 还在有限平面上亚纯，且 \(R\) 包含它的全部有限奇点，留数定理遂给出

\begin{equation}
 \sum_{a\in\C}\operatorname{Res}(f,a)
 +\operatorname{Res}(f,\infty)=0.
 \label{eq:global-residue-sum}
\end{equation}

这里有限奇点确实只有有限多个：它们全在闭圆盘 \(|z|\le R\) 内，而亚纯函数
的极点在有限平面中没有聚点。特别地，有理函数的有限留数之和可以通过无穷远
展开迅速求得。

\section{留数法计算实积分与级数}

留数定理用于实积分时，必须同时证明补上的复路径积分趋于零。下面三种计算
展示不同的轮廓选择。

\subsection{有理函数积分}

若有理函数 \(P/Q\) 在实轴上无极点，且
\(\deg Q\ge\deg P+2\)，取上半圆轮廓。大半圆上
\(P(z)/Q(z)=O(|z|^{-2})\)，弧长为 \(\pi R\)，所以弧积分为 \(O(R^{-1})\)。
令 \(R\to\infty\) 得

\[
 \int_{-\infty}^{\infty}\frac{P(x)}{Q(x)}\,dx
 =2\pi i\sum_{\operatorname{Im}a>0}
 \operatorname{Res}\left(\frac P Q,a\right).
\]

例如 \(1/(z^4+1)\) 在上半平面的极点为
\(e^{i\pi/4}\) 和 \(e^{3i\pi/4}\)，简单极点公式给

\[
 \int_{-\infty}^{\infty}\frac{dx}{x^4+1}
 =\frac{\pi}{\sqrt2}.
\]

\subsection{带分支的积分}

设 \(0<\alpha<1\)。实函数 \(x^{\alpha-1}/(1+x)\) 在零点附近由
\(x^{\alpha-1}\) 控制，在无穷远由 \(x^{\alpha-2}\) 控制，所以两端的反常积分
分别收敛。在 \(0<\arg z<2\pi\) 上取
\(z^{\alpha-1}=e^{(\alpha-1)(\log|z|+i\arg z)}\)。
不能直接把切口当作这一定义域内的积分路径。先固定
\(0<\varepsilon<1/2\)、\(R>2\) 和 \(0<\delta<\pi/2\)，取环扇形
\[
 \varepsilon<|z|<R,\qquad \delta<\arg z<2\pi-\delta
\]
的正向边界：上射线向外，外圆逆时针，下射线向内，内圆顺时针。
轮廓及其邻域避开切口，可对 \(z^{\alpha-1}/(1+z)\) 使用留数定理。

上射线积分为
\[
 \int_\varepsilon^R
 \frac{e^{i\alpha\delta}x^{\alpha-1}}{1+xe^{i\delta}}\,dx;
\]
下射线积分则将 \(\delta\) 换成 \(2\pi-\delta\)，再取负号。
固定 \(\varepsilon,R\) 后令 \(\delta\downarrow0\)，两段被积函数在
\([\varepsilon,R]\) 上一致收敛，径向积分之和趋于

\[
 (1-e^{2\pi i\alpha})
 \int_\varepsilon^R\frac{x^{\alpha-1}}{1+x}\,dx.
\]

两段圆弧也有极限；其积分模分别不超过
\[
 \frac{2\pi R^\alpha}{R-1},\qquad
 \frac{2\pi\varepsilon^\alpha}{1-\varepsilon}.
\]
这些界与 \(\delta\) 无关，随后令 \(R\to\infty\)、\(\varepsilon\downarrow0\)
便均趋于零。轮廓内唯一极点是 \(-1\)，其留数为
\(e^{i\pi(\alpha-1)}\)。故

\[
 \int_0^\infty\frac{x^{\alpha-1}}{1+x}\,dx
 =\frac{2\pi i\,e^{i\pi(\alpha-1)}}{1-e^{2\pi i\alpha}}
 =\frac{\pi}{\sin\pi\alpha}.
\]

分支、切口和两岸相位若有一项未说明，所得常数就可能错误。

\subsection{由亚纯核求和}

函数 \(\pi\cot(\pi z)\) 在每个整数处有留数 \(1\)。对

\[
 F(z)=\frac{\pi\cot(\pi z)}{z^2}
\]

取正方形 \(|\operatorname{Re}z|,|\operatorname{Im}z|\le A_N\) 的正向边界，
其中 \(A_N=N+1/2\)。竖边上
\[
 \cot\bigl(\pi(\pm A_N+iy)\bigr)=-i\tanh(\pi y),
\]
模不超过一。上横边可用
\[
 \cot(\pi z)=i\,\frac{e^{2\pi iz}+1}{e^{2\pi iz}-1}
\]
得到模的上界 \((1+e^{-2\pi A_N})/(1-e^{-2\pi A_N})\le\coth(\pi/2)\)；
下横边由共轭得到同一估计。因此存在与 \(N\) 无关的常数 \(C\)，使
边界上 \(|\cot(\pi z)|\le C\)。又因边界长 \(8A_N\) 且 \(|z|\ge A_N\)，
积分模不超过 \(8\pi C/A_N\)，随 \(N\to\infty\) 趋于零。
非零整数 \(n\) 处留数为 \(1/n^2\)；由

\[
 \pi\cot(\pi z)=\frac1z-\frac{\pi^2z}{3}+O(z^3)
\]

知 \(z=0\) 处留数为 \(-\pi^2/3\)。令 \(N\to\infty\)，留数定理给

\[
 0=-\frac{\pi^2}{3}+2\sum_{n=1}^{\infty}\frac1{n^2},
 \qquad
 \sum_{n=1}^{\infty}\frac1{n^2}=\frac{\pi^2}{6}.
\]

这里级数极限可由绝对收敛控制；不能只写下所有整数留数的形式和而略去轮廓
积分的极限。

\section*{习题}
\addcontentsline{toc}{section}{习题}
\subsection*{A组}
\begin{enumerate}[label=\textbf{\thechapter.\arabic*.}]
\exerciseitem{laurent-geometric-annuli}分别在 \(0<|z|<1\) 与 \(|z|>1\) 展开 \(1/[z(1-z)]\)，并直接核对两个 Laurent 级数的收敛域。
\exerciseitem{laurent-shifted-rational}求 \(1/[(z-1)(z-2)]\) 在以 \(0\) 为中心的三个最大环域中的 Laurent 展开。
\exerciseitem{classify-elementary-singularities}分类 \(z^{-3}\sin z\)、\((1-\cos z)/z^4\)、\(e^{1/z}\) 在 \(0\) 的孤立奇点，并求可求的留数。
\exerciseitem{residue-simple-quotient}设 \(\phi,\psi\) 在 \(a\) 附近全纯，\(\psi(a)=0\)、\(\psi'(a)\ne0\)。从 Taylor 展开证明 \(\operatorname{Res}(\phi/\psi,a)=\phi(a)/\psi'(a)\)。
\end{enumerate}
\subsection*{B组}
\begin{enumerate}[label=\textbf{\thechapter.\arabic*.},resume]
\exerciseitem{removable-small-order}设 \(f\) 在 \(0<|z-a|<R\) 全纯。证明若 \((z-a)f(z)\to0\)，则 \(a\) 可去；并说明仅有 \((z-a)f(z)\) 有界为什么不够。
\exerciseitem{pole-reciprocal-zero}证明 \(f\) 在 \(a\) 有 \(m\) 阶极点，当且仅当 \(1/f\) 在 \(a\) 可去并延拓为 \(m\) 阶零点。
\exerciseitem{residue-higher-pole}证明式\eqref{eq:residue-pole-formula}，并用它计算 \(e^z/(z-1)^4\) 在 \(1\) 的留数。
\exerciseitem{logarithmic-derivative-local-residue}设 \(f\) 在 \(a\) 亚纯且不恒为零。证明 \(f'/f\) 在 \(a\) 的留数等于 \(f\) 的零点阶数减去极点阶数。
\exerciseitem{residue-infinity-coefficient}证明 \(\operatorname{Res}(f,\infty)\) 等于 \(f\) 在无穷远 Laurent 展开中 \(z^{-1}\) 系数的相反数，并计算任意有理函数 \(P/Q\) 的该留数。
\exerciseitem{real-rational-residue-integral}用上半圆轮廓计算
\(\int_{-\infty}^{\infty}(x^2+1)^{-2}\,dx\)，并给出弧积分趋零的明确估计。
\end{enumerate}
\subsection*{C组}
\begin{enumerate}[label=\textbf{\thechapter.\arabic*.},resume]
\exerciseitem{casorati-no-limit}用 Casorati--Weierstrass 定理证明本性奇点处既不可能有有限极限，也不可能有球面意义下的无穷极限。
\exerciseitem{trigonometric-residue-integral}设 \(a>|b|>0\)。令 \(z=e^{it}\)，用留数计算
\(\int_0^{2\pi}(a+b\cos t)^{-1}\,dt\)。
\exerciseitem{cotangent-lattice-sum}对 \(a>0\)，用函数
\(\pi\cot(\pi z)/(z^2+a^2)\) 证明
\[
 \sum_{n=-\infty}^{\infty}\frac1{n^2+a^2}
 =\frac{\pi}{a}\coth(\pi a).
\]
\optionalexerciseitem{keyhole-beta-integral}完整补写钥匙孔轮廓证明，包括分支选取、方向、两段圆弧估计以及常数
\((1-e^{2\pi i\alpha})\) 的来源。
\end{enumerate}
